%------------------------------------------------------------------------------%
% MEASURE AND INTEGRATION THEORY                                               %
%------------------------------------------------------------------------------%
% File  : Measure_and_Integration_Theory.tex                                   %
% Author: Klaus Hoeflinger, Beethovenstrasse 11, D-73117 Wangen                %
% Email : khoeflinger@t-online.de                                              %
%------------------------------------------------------------------------------%

%------------------------------------------------------------------------------%
% DOCUMENT CLASS and INCLUDE FILES                                             %
%------------------------------------------------------------------------------%
\documentclass[11pt,twoside]{article}
\usepackage{geometry}
\geometry{paperheight=241mm,
          paperwidth=152mm,
          top=22mm,
          bottom=17mm,
          left=20mm,
          right=20mm,
          textwidth=112mm,
          textheight=202mm,
          headheight=10mm,
          headsep=3mm,
          footskip=9mm,
          includehead,
          includefoot,
          marginparsep=0mm,
          marginparwidth=0mm}

\renewcommand{\baselinestretch}{0.945}\normalsize

\AtBeginDocument{\setlength\abovedisplayskip{2mm}}
\AtBeginDocument{\setlength\belowdisplayskip{2mm}}

%------------------------------------------------------------------------------%
% define title formats                                                         %
%------------------------------------------------------------------------------%
\usepackage{titlesec}

\renewcommand{\thesection}{\arabic{section}}
\renewcommand{\sfdefault}{FiraSans}
\renewcommand{\ttdefault}{pcr}
\renewcommand{\rmdefault}{antiqua}

\titleformat{\part}[display]
{\normalfont \large \rmfamily \bfseries \centering}
{\small{\thepart.}}{2pt}{}

\titleformat{\section}[runin]
{\normalfont \rmfamily \bfseries}
{\thesection}{1pt}{.\;}[.]

%\titleformat{\section}
%{\normalfont \bfseries \small \filcenter}
%{\thesection.\;}{0mm}{}

\titleformat{\subsection}[runin]
{\normalfont \itshape}
{}{0mm}{}

\titleformat{\subsubsection}[runin]
{\normalfont \itshape}
{}{}{}{}

\titleformat{\paragraph}[runin]
{\normalfont}
{}{}{}{}

\titleformat{\subparagraph}[runin]
{\normalfont}
{}{}{}{}

\titlespacing*{\part}          { 0pt}{ 0pt}{ 8pt}
\titlespacing*{\section}       { 0pt}{ 7pt}{ 4pt}
\titlespacing*{\subsection}    { 0pt}{14pt}{ 6pt}
\titlespacing*{\subsubsection} { 0pt}{ 6pt}{12pt}
\titlespacing*{\paragraph}     { 0pt}{ 6pt}{12pt}
\titlespacing*{\subparagraph}  { 0pt}{ 0pt}{12pt}

%------------------------------------------------------------------------------%
% define enumerate parameters                                                  %
%------------------------------------------------------------------------------%
\usepackage{enumitem}
\setlength{\parindent}{3pt}
\setlength{\parskip}  {0pt}

%\setlist{noitemsep}

\setlist[enumerate]{
         leftmargin=12pt,
         rightmargin=0pt,
         itemsep=1pt,
         itemindent=3pt,
         topsep=3pt,
         labelsep=4pt,
         partopsep=0pt,
         labelwidth=0pt,
         listparindent=0pt,
         parsep=0pt}

\setlist[itemize]{
         leftmargin=12pt,
         rightmargin=0pt,
         itemsep=0pt,
         itemindent=1pt,
         topsep=1pt,
         labelsep=4pt,
         partopsep=1pt,
         labelwidth=0pt,
         listparindent=0pt,
         parsep=0pt}

%------------------------------------------------------------------------------%
% define packages                                                              %
%------------------------------------------------------------------------------%
\usepackage{upref}
\usepackage{amssymb}
\usepackage{slantsc}
\usepackage{amsmath}
\usepackage{amsthm}
\usepackage{thmtools}
\usepackage{amscd}
\usepackage{verbatim}
\usepackage{latexsym}
\usepackage{multicol}
\usepackage{graphicx}
\usepackage{setspace}
\usepackage[ngerman,english]{babel}
\usepackage[protrusion=true,expansion=true]{microtype}
\usepackage{tikz}
\usetikzlibrary{arrows,arrows.meta,decorations.markings}
\usepackage{mathrsfs}
\usepackage{amsfonts}
\usepackage{datetime2}
\usepackage{textgreek}
\usepackage{hyperref}
\usepackage{pifont}
\usepackage{relsize}
%\usepackage{biblatex}
\relsize{-0.05}
\usepackage[skip=0mm]{parskip}
\usetikzlibrary{decorations.pathmorphing}

\usepackage{mlmodern}
\usepackage[T5,T1]{fontenc}
\usepackage{ragged2e}

%\usepackage{mathabx}
%\usepackage{times}

%\usepackage{fontspec}
%\setmainfont{Nimbus Roman No9 L}
%\usepackage[T1]{fontenc}

%------------------------------------------------------------------------------%
% define page layout                                                           %
%------------------------------------------------------------------------------%
\usepackage{fancyhdr}
\pagestyle{fancy}
\setlength{\headheight}{30.0pt}
\renewcommand{\headrulewidth}{0pt}

\AtBeginDocument{
  \addtolength\abovedisplayskip{-0.0\baselineskip}
  \addtolength\belowdisplayskip{-0.0\baselineskip}
}

\fancyhead{}
\fancyfoot{}
\addtolength{\headwidth}{0.02\marginparwidth}

\makeatletter
\let\ps@plain\ps@fancy
\makeatother

\renewcommand\sectionmark[1]
{
 \def\sectionname{#1}
 \markright{\thesection #1}
}

\makeatletter
\@addtoreset{section}{part}
\makeatother

%------------------------------------------------------------------------------%
% define footnote without number                                               %
%------------------------------------------------------------------------------%
\newcommand{\myfootnote}[1]{%
\renewcommand{\thefootnote}{}%
\footnotetext{#1}%
\renewcommand{\thefootnote}{\arabic{footnote}}%
}

%------------------------------------------------------------------------------%
% define numbering in equations                                                %
%------------------------------------------------------------------------------%
\numberwithin{equation}{section}

%------------------------------------------------------------------------------%
% define newtheorem                                                            %
%------------------------------------------------------------------------------%
\declaretheoremstyle[
headfont=\scshape, 
headindent = 0mm, %\parindent,
%postheadhook = {\hspace{0mm}\newline},
%postheadspace = \newline, 
spaceabove = 3mm, 
spacebelow = 3mm,
numberwithin=section,
name=Theorem]{khMATH}
\declaretheorem[style=khMATH]{theorem}

\declaretheoremstyle[
headfont=\bfseries, 
headindent = 0mm, %\parindent,
%postheadhook = {\hspace{0mm}\newline},
%postheadspace = \newline, 
spaceabove = 3mm, 
spacebelow = 3mm,
numberwithin=part,
name=Theorem]{khMATH1}
\declaretheorem[style=khMATH1]{theorem1}

\declaretheoremstyle[
headfont=\scshape, 
headindent = 0mm, %\parindent,
%postheadhook = {\hspace{0mm}\newline},
%postheadspace = \newline, 
spaceabove = 3mm, 
spacebelow = 3mm,
numberwithin=section,
name=Corollary]{khMATH1}
\declaretheorem[style=khMATH1]{corollary}

%            name         counter     label              numeration-mode
\theoremstyle{plain}
\newtheorem*{introduction}            {\sc Introduction}
\newtheorem {standard}                {}                 [section]
\newtheorem*{definition}              {\sc Definition}
\newtheorem{satz1}                    {\sc Satz}
\newtheorem{lemma}                    {\sc Lemma}
\newtheorem*{proposition}             {\sc Proposition}
%\newtheorem{corollary}               {\sc Corollary}
\newtheorem*{corollar}                {\sc Korollar}
\newtheorem*{remark}                  {\sc Remark}
\newtheorem*{remarks}                 {\sc Remarks}
\newtheorem*{example}                 {Example}
\newtheorem*{examples}                {Examples}
\newtheorem*{theo}                    {\sc Theorem}

%------------------------------------------------------------------------------%
% define tabular                                                               %
%------------------------------------------------------------------------------%
\usepackage{tabularx}
\newcolumntype{L}[1]{>{\raggedright\arraybackslash}p{#1}}
\newcolumntype{C}[1]{>{\centering\arraybackslash}  p{#1}} 
\newcolumntype{R}[1]{>{\raggedleft\arraybackslash} p{#1}} 

%------------------------------------------------------------------------------%
% define \sh, \zB                                                              %
%------------------------------------------------------------------------------%
\def\dsh{{d.\,h.}}
\def\ie{{i.\,e.}}
\def\zB{z.\,B.}
\renewcommand{\abstractname}{ }

%------------------------------------------------------------------------------%
% change default spacing for cases                                             %
%------------------------------------------------------------------------------%
\makeatletter
\renewcommand*\env@cases[1][1.6]{%
  \let\@ifnextchar\new@ifnextchar
  \left\lbrace
  \def\arraystretch{#1}%
  \array{@{}l@{\quad}l@{}}%
}
\makeatother

\newenvironment{rcases}
  {\left.\begin{aligned}}
  {\end{aligned}\right\rbrace}

%------------------------------------------------------------------------------%
% change default for \predisplaypenalty                                        %
%------------------------------------------------------------------------------%
\predisplaypenalty=990
\allowdisplaybreaks \numberwithin{equation}{section}

%------------------------------------------------------------------------------%
% general notations                                                            %
%------------------------------------------------------------------------------%
\def\*{\star}
\def\={\,=\,}
\def\+{\,+\,}
\def\xsubset{{\,\subset\,}}
\def\xtimes{{\,\times\,}}
\def\xeof{{\,\in\,}}
\def\eq{{\,=\,}}
\def\xto{{\,\to\,}}
\def\eqdef{\,:=\,}
\def\:{{\,:\,}}
\def\ele{{\,\in\,}}
\def\exists{\mbox{ exists }}
\def\and{\mbox{ and }}
\def\for{\mbox{ for }}
\def\fa{\mbox{ for all }}
\def\fe{\mbox{ for each }}
\def\qfa{\quad \mbox{ for all }}
\def\df{\,\dotfill}
\def\iff{if and only if }
\def\wrt{with respect to}

\makeatletter
\newcommand \Df {\leavevmode \cleaders \hb@xt@ .70em{\hss .\hss }\hfill \kern \z@}
\makeatother

\def\sql{\mathfrak{a}}               % sesquilinear form a
\def\DF{B}                           % Dirichlet form a
%------------------------------------------------------------------------------%
% number sets and special elements                                             %
%------------------------------------------------------------------------------%
\def\P{\mathbb{H}}                   % half plane of $\C$
\def\J{I}                            % interval I
\def\N{{\mathbb{N}}}                 % unsigned integers
\def\Q{{\mathbb{Q}}}                 % rational integers
\def\Z{{\mathbb{Z}}}                 % integers
\def\R{{\mathbb{R}}}                 % real numbers
\def\Rn{{\mathbb{R}}^n}              % real numbers in Rn
\def\Rd{{\mathbb{R}}^d}              % real numbers in Rn
\def\C{{\mathbb{C}}}                 % complex numbers
\def\F{{\mathbb{K}}}                 % arbitrary field
\def\Torus{{\mathbb{T}}}             % complex circle line
\def\T{\tau}                         % upper bound of interval (0,t)
\def\sign#1{\operatorname{sign}(#1)} % sign of a number

%------------------------------------------------------------------------------%
% set theory                                                                   %
%------------------------------------------------------------------------------%
\def\G{G}                            % domain $G$
\def\clG{\overline{\G}}              % closure of $G$
\def\bndG{\partial \G}               % boundary of $G$
\def\setA{\mathcal{A}}               % set A
\def\setB{\mathcal{B}}               % set B
\def\setC{\mathcal{C}}               % set C
\def\setM{M}                         % set M
\def\setN{N}                         % set N
\def\setS{\mathcal{S}}               % set S
\def\famF{\mathfrak{F}}              % family of sets
\def\indI{\mathcal{I}}               % index set I
\def\pot#1{\mathfrak{P}(#1)}         % power set

\def\relR{\mathbf{R}}                % relation R
\def\mapf{f}                         % mapping f
\def\mapg{g}                         % mapping g

\def\set#1#2{\{ \, #1 \,: \, #2 \, \}}
\def\tensor#1#2{#1 \otimes #2}
\def\Tens{\oplus}

\def\cal#1{{\mathcal{#1}}}
\def\aut#1{\operatorname{Aut}(#1)}

\def\cross#1#2{#1 {\,\times\,} #2}
\def\multicross#1#2{#1 \times  \ldots \times  #2}
\def\multitens#1#2 {#1 \otimes \ldots \otimes #2}

%------------------------------------------------------------------------------%
% group theory                                                                 %
%------------------------------------------------------------------------------%
\def\1{{\mathsf{1}}}                % unit element of a monoid
\def\0{{\mathsf{0}}}                % unit element of an Abelian monoid
\def\kernel#1{\mathfrak{N}(#1)}     % kernel of a homomorphism
\def\cokernel#1{\mathfrak{C}{(#1)}} % cokernel of a homomorphism
\def\Hom#1{\operatorname{Hom}(#1)}

%------------------------------------------------------------------------------%
% linear algebra                                                               %
%------------------------------------------------------------------------------%
\def\vecV{\mathit{V}}               % vector space V
\def\vecH{\mathit{H}}               % vector space H
\def\vecW{\mathit{W}}               % vector space W
\def\vecU{\mathit{U}}               % subspace U
\def\convC{\mathcal{C}}             % convex set C
\def\basH{\mathfrak{B}}             % Hamel base B

\def\LO#1{\mathscr{L}(#1) }         % algebra of linear transformations

\def\scalar#1#2{ \langle #1,#2 \rangle }
\def\trace#1{\operatorname{tr}(#1)}
\def\dim#1{{\operatorname{dim} \, #1 }}
\def\codim#1{{\operatorname{codim} \, #1 }}
\def\dim{\operatorname{dim}}
\def\codim{\operatorname{codim}}
\def\ind{\operatorname{ind}}

\def\genG{\cal{G}}

\def\algA{\cal{A}}
\def\algB{\cal{B}}
\def\algU{\cal{U}}

\def\idI{\cal{I}}

%------------------------------------------------------------------------------%
% topology                                                                     %
%------------------------------------------------------------------------------%
\def\compK{{K}}          % compact set C
\def\openO{\mathcal{O}}  % open set O
\def\openU{\mathcal{U}}  % open set U
\def\U{\mathcal{U}}      % neighbourhood U
\def\V{\mathcal{V}}      % neighbourhood V
\def\d{\mathit{d}}       % metric function d

\def\interior#1{\mathring{#1}}
\def\closure#1{{\overline{#1}}}

\def\topT{\mathcal{T}}   % topology T
\def\topX{X}             % topological space X
\def\topY{Y}             % topological space Y
\def\topZ{Z}             % topological space Z
\def\metrS{M}            % metric space M

\def\grpG{G}             % topological group G
\def\grpA{A}             % topological group A
\def\grpB{B}             % topological group B
\def\grpC{C}             % topological group C
\def\grpH{H}             % topological group H
\def\grpU{U}             % topological group U
\def\grpV{V}             % topological group V
\def\topG{G}             % topological group G
\def\topH{H}             % topological group H
\def\topU{U}             % topological subgroup U

\def\subS{\cal{S}}
\def\riR{R}              % ring R

%------------------------------------------------------------------------------%
% calculus                                                                     %
%------------------------------------------------------------------------------%
\def\supp#1{\operatorname{supp}(#1)}
\def\singsupp#1{\operatorname{sing \, supp}(#1)}

\def\re{\operatorname{Re}}
\def\im{\operatorname{Im}}
\def\grad{\operatorname{grad}}

%\def\abs#1 {\left|#1\right|}
\def\abs#1 {\lvert #1 \rvert}
%\def\abs#1 {\left|\,#1\,\right|}

%------------------------------------------------------------------------------%
% measure theory                                                               %
%------------------------------------------------------------------------------%
\def\salgA{\Sigma}               % sigma-algebra S
\def\salgB{\cal{B}}              % sigma-algebra B
\def\Borel#1{\cal{B}_{#1}}       % Borel-algebra 
\def\LB{\lambda}                 % Lebesgue-Borel measure

%------------------------------------------------------------------------------%
% function spaces                                                              %
%------------------------------------------------------------------------------%
\def\Lp{L^p(\G)}                 % spaces L_p
\def\Lq{L^q(\G)}                 % spaces L_q
\def\Lh{L^2(\G)}                 % spaces L_2
\def\Li{L^{\infty}(\G)}          % spaces L_infty
\def\Lc{L_{loc}^1(\G)}           % spaces L_1^{loc}
\def\Ck{C^k(\G)}                 % spaces C_k

\def\BV#1{BV(#1) }               % set of functions of bounded variation
\def\AC#1{AC(#1) }               % set of absolutely continuous functions

\def\MP#1{\mathscr{M}^+(#1) }    % set of positive measures
\def\MS#1{\mathscr{M}^-(#1) }    % vector space of finite signed measures
\def\MC#1{\mathscr{M}(#1) }      % vector space of complex measures
\def\MCR#1{\mathscr{M}_{r}(#1) } % vector space of regular complex measures

\def\SobW{W^{k,p}(\G)}                % Sobolev space W
\def\SobO{\mathring{W}^{k,p}(\G)}     % Sobolev space W
%\def\WN#1#2{\mathring{W}^{#1}(#2) }   % Sobolev space W
\def\WN#1#2{W_0^{#1}(#2) }   % Sobolev space W
%\def\HN#1#2{\mathring{H}^{#1}(#2) }   % Sobolev space H
\def\HN#1#2{H_0^{#1}(#2) }   % Sobolev space H
%------------------------------------------------------------------------------%
% functional analysis                                                          %
%------------------------------------------------------------------------------%
\def\snorm#1{\lVert #1 \rVert}
\newcommand{\norm}[1]{\left\lVert #1 \right\rVert}

\def\topV{V}                       % topological vector space V
\def\topW{W}                       % topological vector space W
\def\normS{X}                      % normed space

\def\banX{\mathit{X}}              % Banach space X
\def\banY{\mathit{Y}}              % Banach space Y
\def\banZ{\mathit{Z}}              % Banach space Z
\def\dbanX{\mathit{X^*}}           % dual space of Banach space X
\def\dbanY{\mathit{Y^*}}           % dual space of Banach space Y
\def\dbanZ{\mathit{Z^*}}           % dual space of Banach space Z

\def\banA{\mathit{A}}              % Banach algebra A
\def\banB{\mathit{B}}              % Banach algebra B
\def\banC{\mathit{C}}              % C*-algebra C


\def\domain#1{\mathfrak{D}(#1)}    % domain
\def\range#1{\mathfrak{R}(#1)}     % range

\def\isoJ{\mathfrak{J}}            % isometry J
\def\repA{{\cal{A}_{\sql}}}        % representation operator A
\def\linS{{S}}                     % linear operator S
\def\linG{{G}}                     % linear operator G
\def\linL{{L}}                     % linear operator L
\def\linT{{T}}                     % linear operator A
\def\linA{\mathit{A}}              % linear operator A
\def\linB{{B}}                     % linear operator B
\def\linM{{M}}                     % linear operator M
\def\linU{{U}}                     % linear operator U
\def\linI{{I}}                     % linear operator I
\def\A{\cal{A}}                    % linear differential operator A
\def\BDO{\cal{B}}                  % linear boundary differential operator B
\def\linU{{U}}                     % unitary operator U
\def\linP{{P}}                     % projection P
\def\compA{k}                      % compact operator k
\def\linFR{f}                      % finite rank operator f
\def\fredA{A}                      % Fredholm operator F

\def\HAM{\cal{H}}                  % Hamilton operator
\def\PA{\partial^{\alpha}}         % elementary differential operator
\def\DO{\cal{A}}                   % linear differential operator
\def\BC{\cal{B}}                   % boundary operator
\def\SLO{\cal{L}_{p,q}}            % Sturm-Liouville operator
\def\MO{\cal{L}_{M}}               % momentum operator

\def\HS#1{S(#1)           }        % algebra of Hilbert-Schmidt operators
\def\FR#1{F(#1)           }        % algebra of finite rank operators
\def\TC#1{N(#1)           }        % algebra of trace class operators
\def\BU#1{\mathscr{U}(#1) }        % algebra of unitary linear operators
\def\BO#1{\mathscr{L}(#1) }        % algebra of bounded linear operators
\def\CO#1{\mathscr{C}(#1) }        % algebra of compact linear operators
\def\CL#1{\mathscr{C}(#1) }        % algebra of compact linear operators
\def\FO#1{\Phi(#1) }               % set of Fredholm operators
\def\FK#1#2{\Phi_{#2}(#1) }        % set of Fredholm operators with index k

\def\hilH{\mathit{H}}              % Hilbert space H
\def\clU{\mathit{U}}               % closed subspace U

\def\orthO{\mathfrak{O}}           % orthonormal system O
\def\orthS{\mathfrak{S}}           % orthonormal system S
\def\basB{\mathfrak{B}}            % basis B of a Hilbert space

\def\GL#1 {\mathbf{GL}(#1)}        % linear group
\def\OG#1 {\mathbf{O}(#1)}         % linear group
\def\SOG#1 {\mathbf{SO}(#1)}       % linear group

%------------------------------------------------------------------------------%
% distribution theory                                                          %
%------------------------------------------------------------------------------%
\def\disF{F}                  % distribution F
\def\disG{G}                  % distribution G
\def\disH{H}                  % distribution H
\def\disU{U}                  % distribution U
\def\disT{T}                  % distribution T
\def\SF{C^{\infty}(\G)}       % space of smooth functions
\def\TF{C_c^{\infty}(\G)}     % space of compactly supported smooth functions
\def\TFRd{C_c^{\infty}(\R^d)} % space of test functions
\def\SRd{\mathscr{S}(\R^d)}   % Schwartz space
\def\SR1{\mathscr{S}(\R)}     % Schwartz space
\def\B1#1{C^1(#1) }           % space of continuously differentiable functions
\def\Bm#1{C^m(#1) }           % space of continuously differentiable functions
\def\Ck{C^k(\G)}
\def\TD{\mathscr{S}'(\R^d)}   % space of temperated distributions
\def\E{\cal{E}}               % space of distributions with compact support
\def\D{\mathscr{D}}           % D
\def\FT{\mathcal{F}}          % Fourier transformation F
\def\FT{\mathscr{F}}          % Fourier transformation F

%------------------------------------------------------------------------------%
% other definitions                                                            %
%------------------------------------------------------------------------------%
\def\Proof{\textsc{Proof}}
\def\FRECHET{Fr\'{e}chet}
\def\ARZELA{Arzel\'{a}}
\def\GARDING{G\r{a}rding}
\def\HOELDER{H\"older}
\def\SCHROEDINGER{Schr\"odinger}
\def\JOERGENS{J\"orgens}
\def\WUEST{W\"uest}
\def\KAEHLER{K\"aehler}
\def\HOEFLINGER{H\"oflinger}
\def\POINCARE{Poincar\'{e}}
\def\FA{Functional Analysis}

\def\Author   {K.\,{\HOEFLINGER}}
\def\AuthorUC {K.\,HOEFLINGER}
\def\Version  {1.0}
\def\Email    {khoeflinger@t-online.de}


\titleformat{\part}%[display]
{\normalfont \rmfamily \bfseries \centering}
{{\thepart.}}{3pt}{}

\titleformat{\section}[runin]
{\normalfont \bfseries \filcenter}
{\thesection.\,}{0mm}{}[.]

\titlespacing*{\part}   {5pt}{12pt}{ 4pt}
\titlespacing*{\section}{0pt}{ 3pt}{ 3pt}

%------------------------------------------------------------------------------%
%                                BEGIN OF DOCUMENT                             %
%------------------------------------------------------------------------------%
\begin{document}
\thispagestyle{empty}

\centerline{\textrm{\large{\textbf{Measure and Integration Theory}}}}
\vspace{4pt}
\centerline{\textsc{\small{\Author}}}
\vspace{10pt}

\fancyhead[OL]{\small{\thepage}}
\fancyhead[ER]{\small{\thepage}}

%------------------------------------------------------------------------------%
%                                  CONTENT                                     %
%------------------------------------------------------------------------------%
\myfootnote
{
  Version {\Version} from \today;
  Email:  \textit{\Email}
}

\thispagestyle{empty}
This text serves as a comprised review on the essentials
of \textit{measure and integration theory},
a fundamental area in modern mathematics with many applications
to functional and harmonic analysis and other fields of importance.
The content is presented within three parts:
the first one describes abstract measure spaces and different kinds of
measures, the second one deals with \textsc{Lebesgue}'s integration theory
and the final one considers measures on locally compact topological spaces.

\fancyhead[OC] {\small{\textsc{I. Measure Spaces}}}
\fancyhead[EC] {\small{\textsc{I. Measure Spaces}}}
\part{Measure Spaces}
%-------------------------------------------------------------------------------
In mathematics a \textit{measure} is
a way of assigning a numerical value to certain subsets of a given base set,
which can be considered as the length, the area or the volume of such sets.
One may think of the situation in the real line,
where each compact interval $[a,b]$ is measured by the real number $b{-}a$.
This mapping can easily be extended to finite or actually countable
many unions of finite intervals under the assumption
that the involved intervals are pairwise disjoint.

\section{Measurable Spaces}
%------------------------------------------------------------------------------%
Motivated by the introductory statements,
we are interested in collections of subsets of a set $X$,
which are closed under taking countable unions,
and each member of such a collection will be called a \textit{measurable set}.
This means
that the collection of measurable sets will be the natural domain for a measure.

\subparagraph{}
For a non-empty set $X$,
the domain of a measure on $X$ is a collection of subsets of $X$,
called a \textit{$\sigma$-algebra}.
In particular,
a $\sigma$-algebra $\salgB$ on $X$ is defined by the following properties:

\begin{itemize}[leftmargin=15pt]
\item[1.]
$\salgB$ contains the empty set $\emptyset$.

\item[2.]
$\salgB$ is closed under complemention:
if $A,B \xeof \salgB$, then $A \setminus B \xeof \salgB$
(this property implies $X \xeof \salgB$).

\item[3.]
$\salgB$ is closed under taking countable unions,
that is,
if $\{ {A_n}\}_{n=1}^{\infty}$ is a sequence of sets $A_n \xeof \salgB$,
then $\bigcup_{n=1}^{\infty} A_n \xeof \salgB$.
\end{itemize}

\subparagraph{}
From de Morgan's law
(which is just as valid for infinite unions and intersections
as it is for finite ones),
one can see
that $\sigma$-algebras are also closed under countable intersections.
We call an element of a $\sigma$-algebra a \textit{measurable set},
and each pair $(X,\salgB)$ consisting of a set $X$ and
a $\sigma$-algebra $\salgB$ is said to be a \textit{measurable space}.

\section{Generators of \textsigma-Algebras}
%------------------------------------------------------------------------------%
We now describe two methods,
how $\sigma$-algebras can be established supposed a base set $X$ is given.

\subparagraph{}
First,
we may claim
that the intersection of arbitrarily many $\sigma$-algebras
forms a $\sigma$-algebra again.
Thus,
for a given collection $\setA$ of subsets of the set $X$
(think of the collection of finite intervals in $\R$ for instance),
we define the {coarsest} $\sigma$-algebra containing $\setA$
(denoted $[\setA]$).
This resulting $\sigma$-algebra $[\setA]$ is called
\textit{generated} by $\setA$,
and $\setA$ is said to be a \textit{generator} of $[\setA]$.

\subparagraph{}
Second,
if $X'$ is a subset of $X$ and $\salgB$ is a $\sigma$-algebra on $X$,
it is possible to define the \textit{subset-$\sigma$-algebra} on $X'$
as the collection of all intersections $X' \cap B$ with $B \xeof \salgB$.

\paragraph*{}
Using the first method described above,
we define the \textit{Borel} $\sigma$-algebra $\cal{B}(\R)$
on the real line $\R$ to be the $\sigma$-algebra generated
by the collection of closed and finite intervals $[a,b] \subset \R$.
Each element of $\cal{B}(\R)$ is called a \textit{Borel-measurable set}.
Similarly,
for the higher dimensional Euclidean spaces
the Borel $\sigma$-algebras $\cal{B}(\Rn)$ are 
generated by the collection of open subsets, closed subsets,
compact subsets and open balls in $\Rn$.
Applying the second method for any subset $X \subset \Rn$,
we can define the $\sigma$-algebra $\cal{B}(X)$,
which is also called the Borel $\sigma$-algebra on $X$.

\section{Positive Measures}
%------------------------------------------------------------------------------%
Given any collection of measurable sets,
a measure is nothing but a \textit{set function} acting additively
on sequences of pairwise disjoint measurable sets.

\subparagraph{}
Let $(X,\salgB)$ a measurable space.
Then a \textit{measure} is a mapping
$\mu \: \salgB \xto \R^+ \cup \{\infty\}$
satisfying (1) $\mu(\emptyset)=0$ and (2) the main condition
\[
\mu (\, \bigsqcup_{i=1}^{\infty} M_i \,) \= \sum_{i=1}^{\infty} \, \mu(M_i)
\]
for any sequence of pairwise disjoint sets
$\{M_1,M_2,M_3, \ldots\} \subset \salgB$,
which is the fundamental property of a measure
denoted \textit{countable additivity} or \textit{$\sigma$-additivity}.

\subparagraph{}
Since a measure $\mu$ takes values in the interval $[0,\infty]$,
it is often called an \textit{unsigned} or a \textit{positive} measure
(in distinction to other concepts like 
\textit{signed} and \textit{complex} measures which are discussed below).

\subparagraph{}
Each triple $(X, \salgB, \mu)$ consisting of a measurable space
and a measure $\mu$ on $X$ is called a \textit{measure space}.
It is a central concept in measure theory,

\paragraph{}
Given a measure $\mu$,
the inclusion $A \subset B$ implies $\mu(A) \leq \mu(B)$
for any measurable sets $A,B$.
This property of a measure is called \textit{monotonicity}.
But there are further properties which measures may or may not have:

\begin{itemize}[leftmargin=12pt]
\item[1.]
If $\mu(A) < \infty$ for each set $A \xeof \salgB$,
then $\mu$ is called a \textit{finite} measure.
The measure $\mu$ is called $\sigma$-finite,
if there is a sequence ${A_n}$ of measurable sets
such that
\[
X = \bigcup_{n \in \N} A_n, \quad A_n \subset A_{n+1}, \quad
\mu(A_n) < \infty \, \fa n \in \N.
\]

\item[2.]
A measure $\mu$ is said to be \textit{concentrated} on a measurable set $A$,
if $\mu(E) \eq 0$ for all sets $E \xeof \salgB$ which are distinct from $A$.
This is equivalent to the statement
that $\mu(E) \eq \mu(A \cap E)$ for every $E \xeof \salgB$.

\item[3.]
For two measures $\mu,\nu$ on $\salgB$
suppose there exist disjoint sets $A_{\mu},A_{\nu}$ such that
$\mu$ is concentrated on $A_{\mu}$ and
$\nu$ is concentrated on $A_{\nu}$.
Then $\mu, \nu$ are called \textit{mutually singular} denoted $\mu \bot \nu$.
\end{itemize}

\subparagraph{}
Finally,
a property $E$ defined for elements $x \xeof X$ is called
\textit{$\mu$-almost everywhere valid},
if the set of all elements $x$,
for which $E(x)$ does \textit{not} hold,
is a null set or the subset of another null set.

\section{Signed Measures}
%------------------------------------------------------------------------------%
There are several possibilities to extend the notion of a (positive) measure.
One of them is that of a \textit{signed measure},
which is characterized by the property
to take values in the extended real lines
$(-\infty,\infty] := \R \cup \{\infty\}$
or
$[-\infty,\infty) := \R \cup \{-\infty\}$.
In particular,
for a measurable space $(X,\salgB)$ a set function
\[
\mu: \salgB \to (-\infty,+ \infty] \quad \mbox{or} \quad
\mu: \salgB \to [-\infty,+ \infty)
\]
is called a \textit{signed measure},
if $\mu(\emptyset) \eq 0$ and $\mu$ is countable additive.
A signed measure is called \textit{finite},
if it takes values in the real line only.

\subparagraph{}
Each positive measure is a special case of a signed measure.
If $\mu,\nu$ are signed measures,
one can take the difference $\mu - \nu$, supposed one of them is finite.
Conversely,
for signed measures the \textsc{Hahn-Jordan} \textit{decomposition theorem}
holds,
which states
that each signed measure can be represented as the difference of two
positive measures:

\begin{theorem}\label{HAHN_JORDAN_DECOMPOSITION_THEOREM}
For each signed measure $\mu$ on the measurable space $(X,\salgB)$
there are two positive measures $\mu_+$ and $\mu_-$,
such that for every set $A \xeof \salgB$ the equality
$\mu(A) \eq \mu_+(A) - \mu_-(A)$ holds.

\subparagraph{}
In addition, 
there is a disjoint decomposition of the base set $X$ into two sets $P,N$,
such that $\mu_+$ is the zero measure on $N$
and $\mu_-$ is the zero measure on $P$.
If $X \eq P' \sqcup N'$ is a further decomposition with these properties,
then the sets $P,P'$ and $N,N'$ differ only in null sets {\wrt} $\mu$.
\end{theorem}
  
\subparagraph{}
The sum of two finite signed measures is a finite signed measure
as well as the product of a finite signed measure by a real number.
Hence finite signed measures are closed under linear combinations of them,
which implies
that the set of finite signed measures on a measurable space $(X,\salgB)$
is a real linear space denoted $\cal{M}^-(X,\salgB)$.
For a signed measure $\mu$,
the measure $\abs{\mu} := \mu_+ + \mu_-$ is called
the \textit{total variation} of $\mu$,
which itself is called of \textit{bounded variation}, if $\abs{\mu} < \infty$.
 
\section{Complex Measures}
%------------------------------------------------------------------------------%
In contrast to positive measures,
a \textit{complex measure} on the measurable space $(X,\salgB)$
is a $\sigma$-additive set function $\mu \: \salgB \xto \C$.
By this reason,
a complex measure can only have finite values,
and the condition $\mu(\emptyset)=0$ is fulfilled by itself
and need not be part of the definition.

\subparagraph{}
When dealing with complex measures,
it is possible to define the concept of \textit{variation}.
More precisely,
the variation $\abs{\mu} $ of a complex measure $\mu$
is that mapping $\abs{\mu} \: \salgB \xto [0,\infty]$,
which assigns to each measurable set $A$ the value
\[
\abs{\mu} (A) \eqdef
\sup \, \sum_{n=1}^{\infty} \, \norm{\mu(A_n)},
\]
where the sequences $(A_n)$ may be any mostly countable disjoint
decompositions of $A$ in measurable subsets.
It is a surprising result
that for each complex measure $\mu$
the variation $\abs{\mu} $ is actually a \textit{finite} positive measure.

\subparagraph{}
The \textit{total variation} of a complex measure $\mu$ is defined
to be the variation of the entire set $X$,
that is, $\abs{\mu} := \abs{\mu} (X)$.
Since $\abs{\mu} < \infty$,
each complex measure is said to be of \textit{bounded variation}.

\section{Measurable Functions}
%------------------------------------------------------------------------------%
The main thing what one can do with a measure is to integrate
real-valued or complex-valued functions.
But in general this is not possible for any function, that is,
the restriction to the so-called \textit{measurable functions} is necessary.
By this way
we will establish the \textit{Lebesgue integral} for measurable functions,
which may or may not exist for such a function.
It turns out
that the Lebesuge integral has much better behaviour {\wrt}
some limit processes than the classical Riemann integral,
which suffers from the drawback
that the limit function of a sequence of Riemann-integrable functions
need not be Riemann-integrable again.

\paragraph{}
The notion of measurability of a function is elementary
in integration theory.
Let $(X,\salgB)$ and $(X',\salgB')$ be measurable spaces,
then any mapping $f \: X \xto X'$ is called \textit{measurable},
if $f^{-1}(A) \xeof \salgB$ for every set $A \xeof \salgB'$.
Note
that this definition is rather similar to that of continuity
for functions between topological spaces.

\subparagraph{}
The measurablility of a given function $f$ is ensured,
if the pre-images of those subsets in $X'$ are measurable in $X$,
which form a generator of the $\sigma$-algebra $\salgB'$.
Especially in case of a measurable space $(X,\salgB)$
and a real-valued function $f \: X \xto \R$,
the function $f$ is measurable {\wrt} the $\sigma$-algebra
$\cal{B}(\R)$, if and only if the pre-images of the half-bounded intervals
$A_{\lambda} := \set{x \xeof \R}{x \leq \lambda}$ are elements of $\salgB$,
and in this case the mapping $f$ is called a \textit{Borel function}.

\subparagraph{}
If $f,g$ are Borel functions and $\lambda$ is a real number,
then $f+g$ and $\lambda f$ are Borel functions again,
so the set of Borel functions for a given measurable space
forms a linear space.
It is also important to know that $f$ is measurable,
iff the functions $f^+$ and $f^-$ are both measurable,
which are defined to be the positive and negative part
of the real-valued function $f$,
that is, $f \eq f^+ - f^-$.
We conclude
that the modulus function $\abs{f} := f^+ + f^-$ is measurable,
supposed the function $f$ itself is measurable.

\subparagraph{}
In case of complex-valued functions,
a function $f \: X \xto \C$ is said to be measurable,
iff both the real and imaginary part of $f$ are measurable.

%------------------------------------------------------------------------------%
%                                 BIBLIOGRAPHY                                 %
%------------------------------------------------------------------------------%
\begin{comment}
\vspace{8mm}
\begin{small}
\begin{thebibliography}{99}

\bibitem{El_1} J.\,Elstrodt:
\textit{Ma{\ss} -und Integrationstheorie}.
Springer Verlag, 2018.

\end{thebibliography}
\end{small}
\end{comment}

%\newpage
%\thispagestyle{empty}
\fancyhead[OC]{\small{\textsc{II. Lebesgue Integration}}}
\fancyhead[EC]{\small{\textsc{II. Lebesgue Integration}}}
\part{Lebesgue Integration}
%-------------------------------------------------------------------------------
Integration is a way of assigning a mean value to a given function.
The basic idea of \textit{Lebesgue integration} is to decomposite
the range of a function $f \: X \xto \R$ into intervals
and to measure the pre-images of them.
If these pre-images are measurable sets {\wrt} a measurable space $(X,\salgB)$,
we can do a first step to assign a real number to $f$,
which may be regarded as the integral of $f$.

\section{Setup of the Lebesgue Integral}
%------------------------------------------------------------------------------%
Usually the definition of the Lebesgue integral is carried out in four steps,
beginning with that for elementary functions.
Let $(X, \salgB, \mu)$ be a measure space.
Our aim is to define Lebesgue-integration for every measurable function,
which will be done by a couple of steps:

\begin{itemize}[leftmargin=12pt,itemsep=2pt]
\item[1.]
A mapping $\chi_A \: X \xto \R$ is called the \textit{characteristic function}
of a measurable set $A$,
if $\chi(x)=1$ for $x \xeof A$ and $\chi(x)=0$ otherwise.
We can simply define the Lebesgue-Integral of $\chi_A$ by
\[
\int_{X} \chi_A \, d \mu \eqdef \mu(A),
\]
and the function $\chi_A$ is called \textit{Lebesgue-integrable},
iff $\mu(A) < \infty$.

\item[2.]
A measurable function $e \: X \xto \R$ is called \textit{simple},
if there are finitely many real numbers $a_1,\ldots,a_n$
such that $e(x) \xeof \{a_1,\ldots,a_n\}$ for all $x \xeof \R$,
that is, the range of a simple function is a finite set.
Equivalently,
a simple function $e$ is a linear combination of characteristic functions,
and we define the Lebesgue integral for $e$ by the setting
\[
\int_{X} e \, d \mu \eqdef
\sum_{i=1}^m a_i \cdot \int_{X} \chi_{A_i} \, d \mu.
\]
If this is a finite value,
then the simple function $e$ is said to be \textit{Lebesgue-integrable}.

\item[3.]
Now,
let the function $f \: X \xto \R$ be nonnegative and measurable.
In this case we define the set $\cal{E}_f$ as the collection
of all simple functions $e$ satisfying the property $e(x) \leq f(x)$
for all $x \xeof X$.
Note
that the set $\cal{E}_f$ is non-empty,
since for each measurable function $f$
there is a sequence $\{e_n\}$ of simple functions with
$e_1 \leq e_2 \leq e_3 \leq \ldots$ and
$\lim_{n \to \infty} e_n \eq f$ pointwise.
We define the \textit{Lebesgue integral} of $f$ as follows:
\[
\int_{X} f \, d \mu :=
\sup_{e \in \cal{E}_f} \, \int_{X} e \, d \mu.
\]
If the supremum is a finite number,
the nonnegative and measurable function $f$ is called
\textit{Lebesgue-integrable}.

\item[4.]
In order to define the Lebesgue integral
of a measurable function $g \: X \xto \R$,
let us decomposite $g$ into the difference of nonnegative measurable functions
$f^+, f^-$, i.e. $g := f^+ - f^-$,
reducing the concept of integration for $g$
to that of the functions $f^+$ and $f^-$.
Then the function $g$ is said to be \textit{Lebesgue-integrable},
if $f^+$ and $f^-$ are both Lebesgue-integrable,
and the Lebesgue integral of $g$ is defined according to
\[
\int_{X} g \, d \mu \eqdef
\int_{X} f^+ \, d \mu \, - \, \int_{X} f^- \, d \mu,
\]
It turns out that a measurable function $g$ is Lebesgue-integrable,
if and only if the function $\abs{g} := f^+ + f^-$ is Lebesgue-integrable,
and in this case the function $f$ is said to be \textit{absolutely integrable}.
\end{itemize}

\subparagraph{}
The construction of the Lebesgue integral is done now.
Lebesgue integration enjoys a list of nice properties:

\begin{itemize}[leftmargin=12pt]
\item[1.]
If $A \subset X$ is a null set,
the Lebesgue integral has the value $0$ for any measurable function.

\item[2.]
If two functions $f,g$ are Lebesgue-integrable,
then any linear combination $\alpha f + \beta g$ with $\alpha, \beta \xeof \R$
is Lebesgue-integrable, too, and the equation
\[
\int_X (\alpha f + \beta g) \, d \mu =
\alpha \int_X  f \, d \mu + \beta \int_{X}  g \, d \mu
\]
holds.
Thus the mapping $f \mapsto \int_{X} f \, d \mu$ is a linear map.
More precisely, it is  a linear functional
on the $\R$-linear space of Lebesgue-integrable functions.

\item[3.]
For a Lebesgue-integrable function the following inequality holds:
\[
\abs{\int_X f \, d\mu \,} \, \leq \, \int_{X} \abs{f} \, d\mu \,
\]

\item[4.]
If $f \leq g$ $\mu$-almost everywhere, then
\[
\int_X f \, d\mu \, \leq \, \int_{X} g \, d\mu,
\]
\end{itemize}

\subparagraph{}
Extensions of the Lebesgue integral are possible to
(1) complex-valued measurable functions and
(2) measurable subsets of the base space $X$.
Firstly,
for the class of complex-valued functions
we define $f$ to be Lebesgue-integrable,
iff real and imaginary part of $f$ are so,
and in this case the Lebesgue integral is
\[
\int_{X} f \, d \mu
\eqdef
\int_{X} \re f \, d \mu
\, + \, i \cdot
\int_{X} \im f \, d \mu.
\]
The mapping $f \mapsto \int f d\mu$ is $\C$-linear on the linear space
of Lebesgue-integrable complex-valued functions.
As in case of real-valued functions,
a Lebesgue-integrable complex-valued function satisfies the inequality
\[
\abs{\int_{X}      f  \, d\mu \,} \, \leq \,
     \int_{X} \abs{f} \, d\mu.
\]
Secondly,
if $A \xeof \salgB$ is measurable, then
\[
\int_{A} f \, d\mu \eqdef \int_{X} f \cdot \chi_A \, d\mu.
\]

\subparagraph{}
For a fixed measure space $(X,\salgB,\mu)$,
one can construct a lot of further measures using Lebesgue integration
by the setting
\[
\mu_f(A) \eqdef \int_X f \cdot \chi_A \, d \mu, \quad A \in \salgB,
\]
where $f \: X \xto [0,\infty]$ is any nonnegative measurable function.
If $g$ is another function fullfilling these properties,
the formula
\[
\int_X g \, d \mu_f \= \int_X g \cdot f \, d \mu
\]
holds,
and if $\mu$ is $\sigma$-finite, the identity $\mu_f \eq \mu_g$ is valid,
if and only if
$f \eq g$ $\mu$-nearly everywhere.
Thus a function $f \: X \xto [0,\infty]$
(modulo $\mu$-almost everywhere equivalence)
can be identified with a positive measure $\mu_f$ on $X$.

\section{Convergence Theorems}
%------------------------------------------------------------------------------%
One of the main purposes of the Lebesgue integral is,
in contrary to the classical Riemann integral,
that under some restrictions
we can swap the limit process of a pointwise convergent sequence
of functions with integration in the sense that
\[
\lim_{n \to \infty} \int_{X} f_n \, d \mu
\=
\int_{X} \lim_{n \to \infty} f_n \, d \mu.
\]
There are some cases in which this exchange can be done without doubt,
and we will mention two theorems which describe such a desirable situation.

\paragraph{}
The first one,
due to \textsc{B.\,Levi},
is called the \textit{monotone convergence theorem} and states as follows:

\begin{theorem}\label{MONOTON_CONVERGENCE_THEOREM}
Let the real-valued functions $\{f_n\}_{n=1}^{\infty}$ be nonnegative,
Lebesgue-integrable and monotone increasing, i.e.
\[
f_i(x) \leq f_j(x) \, \mbox{ for all } x \in X, \quad i < j.
\]
Then there is a Lebesgue-integrable function $f$ having the property
that $f \eq \lim_{n \to \infty} f_n$ $\mu$-almost-everywhere and
\[
\lim_{n \to \infty} \int_{X} f_n \, d \mu \,
= \,
\int_{X} \lim_{n \to \infty} f_n \, d \mu.
\]
\end{theorem}

\subparagraph{}
Thus the process of limitation does not exceed
the class of Lebesgue-integrable functions.
In this sense the set of Lebesgue-integrable functions is closed.

\paragraph{}
The second theorem,
called the \textit{dominated convergence theorem} due to \textsc{H.\,Lebesgue},
yields a sufficient criterion for integrability of a function,
which is defined by the limit process of a sequence of integrable functions:

\begin{theorem}\label{DOMINATED_CONVERGENCE_THEOREM}
Suppose that the sequence $(f_n)$ of Lebesgue-integrable functions
$f_n \: X \xto \R$ has the property
\[
\lim_{n \to \infty} f_n = f \quad \mu \mbox{-almost everywhere}.
\]
In addition, 
if there is a Lebesgue-integrable function $g$
with $\abs{f_n} \leq g$ for all $n \xeof \N$,
then $f$ is Lebesgue-integrable, too,
and integration can be commuted with the process of limitation,
that is
\[
\int_{X} f \, d \mu\,
= \,
\lim_{n \to \infty} \int_{X} f_n \, d \mu.
\]
\end{theorem}

\section{Product Spaces}
%------------------------------------------------------------------------------%
The method of constructing a $\sigma$-algebra by given subsets
of a base set can be used to define a $\sigma$-algebra over the
cartesian product of measurable spaces.

\subparagraph{}
Let $(X_1,{\salgB}_1)$ and $(X_2,{\salgB}_2)$ be measurable spaces.
Then for the cartesian product $\cross{X_1}{X_2}$
there is a $\sigma$-algebra $\cross{{\salgB}_1}{{\salgB}_2}$
called the \textit{product $\sigma$-algebra},
which is generated by all sets $\cross{A_1}{A_2}$
with $A_1 \xeof \salgB_1, A_2 \xeof \salgB_2$.
It is the smallest $\sigma$-algebra over $\cross{X_1}{X_2}$,
such that the projection mappings $(x_1,x_2) \mapsto x_1$ and
$(x_1,x_2) \mapsto x_2$ are still measurable.
This construction can be extended to a finite number of products
of measurable spaces.

\subparagraph{}
If $\mu_1,\mu_2$ are measures on the measurable spaces
$(X_1,{\salgB}_1)$ and $(X_2,{\salgB}_2)$,
then there is a measure $\mu_1 \otimes \mu_2$ on the $\sigma$-algebra
$\cross{{\salgB}_1}{{\salgB}_2}$,
such that for all measurable sets $A_1 \xeof \salgB_1$, $A_2 \xeof \salgB_2$:
\[
\mu_1 \otimes \mu_2(A_1,A_2) \= \mu_1(A_1) \cdot \mu_2(A_2)
\]
The measure $\mu_1 \otimes \mu_2$ obtained this way is called
the \textit{product measure} on the measurable space
$(\cross{X_1}{X_2}, \cross{\salgB_1}{\salgB_2})$.
If both measures $\mu_1,\mu_2$ are $\sigma$-finite,
then $\mu_1 \otimes \mu_2$ is unique and $\sigma$-finite, too.

\paragraph{}
Under certain conditions
the Lebesgue integral for any measurable function
$f \: X_1 \otimes X_2 \xto \R$ {\wrt} the product measure $\mu_1 \otimes \mu_2$
can be reduced to the Lebesgue integrals of the single measures.
This is the assertion of the \textsc{Fubini-Tonelli} \textit{theorem}:

\begin{theorem}
Let $f \: X_1 \xtimes X_2 \xto \R$ be $\mu_1 \otimes \mu_2$-integrable.
Then the following holds:

\begin{itemize}[leftmargin=12pt]
\item[1.]
the function $f_y(x) := f(x,y)$ is $\mu_1$-almost-everywhere integrable,

\item[2.]
the function $f_x(y) := f(x,y)$ is $\mu_2$-almost-everywhere integrable,

\item[3.]
the order of integration can be commuted, that is, if
\[
M_1(x) \eqdef \int_{X_2} f_x(y) \, d \mu_2,
\quad \quad
M_2(y) \eqdef \int_{X_1} f_y(x) \, d \mu_1,
\]
then the following identities hold:
\[
\int_{X_1 \xtimes X_2}
f(x,y) \, d(\mu_1 \otimes \mu_2)
\=
 \int_{X_1}
M_1(x) \, d \mu_1
\=
 \int_{X_2}
M_2(y) \, d \mu_2.
\]
\end{itemize}
\end{theorem}

\section{Absolutely Continuous Measures}
%------------------------------------------------------------------------------%
The concept of an \textit{absolutely continuous measure} is of 
fundamental meaning in measure theory.
Let $(X,\salgB,\mu)$ be a measure space
and $\nu$ a further measure define on the measurable space $(X,\salgB)$.
Then $\nu$ is called \textit{absolutely continuous} {\wrt} $\mu$
(notation: $\nu \ll \mu$),
if each $\mu$-null set is always a $\nu$-null set,
that is, $\mu(A) \eq 0$ implies $\nu(A) \eq 0$ for each measurable set $A$.

\subparagraph{}
Two measures are said to be equivalent,
if each is absolutely continuous {\wrt} the other one.
Hence we can divide all the measures defined on a measurable space
into equivalence classes and denote each of them a \textit{measure class}.

\subparagraph{}
Given any positive measure $\mu$,
the related absolutely continuous measures {\wrt} $\mu$ can easily be found.
If $f \: X \xto [0,\infty)$ is a Lebesgue integrable function,
then $f$ causes by means of the mapping
\[
\nu_f(A) \=
\int_{A} f \, d\mu, \quad A \in \salgB
\]
an absolutely continuous measure $\nu_f$ {\wrt} $\mu$.
Thus Lebesgue integrable functions can be regardes as
absolutely continuous measures on the measurable space $(X,\salgB)$.
On the other side,
the \textit{Radon-Nikodym} \textit{theorem} states that
the measure $\nu_f$ are the only absolutely continuous measures {\wrt} $\mu$:

\begin{theorem}\label{RADON_NIKODYM_THEOREM_A}
Let $\mu$ be a $\sigma$-finite measure on $(X,\salgB)$,
$(X,\salgB,\mu)$ be a $\sigma$-finite measure space
and $\nu$ a measure being absolutely continuous {\wrt} $\mu$.
Then there exists a Lebesgue-integrable map $f_{\nu} \: X \xto [0,\infty)$
such that
\[
\nu(A) \= \int_{X} f_{\nu}\, d \, \mu
\quad \mbox{for all } A \in \salgB.
\]
\end{theorem}

\subparagraph{}
The condition of $\mu$ being $\sigma$-finite is essential.
The function $f_{\nu}$ is called the \textit{Radon-Nikodym derivation}
of $\nu$ {\wrt} $\mu$ and is uniquely defined 
$\mu$-mostly everywhere.

\subparagraph{}
The {Radon-Nikodym theorem} can also be stated
for a complex measure $\nu$,
but one has to use the variation $\abs{\nu} $ in this case,
which is known to be a positive finite measure.
A complex measure $\nu$ is called \textit{absolutely continuous}
{\wrt} the positive measure $\mu$,
if $\abs{\nu} \ll \mu$.
Then a further variant of the {Radon-Nikodym theorem} reads as follows:

\begin{theorem}\label{RADON_NIKODYM_THEOREM_B}
Let $\mu$ be a $\sigma$-finite positive measure on the measurable space
$(X,\salgB)$ and $\nu$ be a complex measure satisfying $\nu \ll \mu$.
Then there is $f \xeof L^1(X;\C)$
such that
\[
\nu(A) \= \int_{X} f \, d\mu
\quad \mbox{for all } A \in \salgB.
\]
\end{theorem}

\paragraph{}
We will finish the considerations about absolutely continuous measures
with the \textsc{Lebesgue} \textit{decomposition theorem}.
It asserts
that for a given $\sigma$-finite measure space $(X,\salgB,\mu)$
each complex measure $\nu$ can be represented as the sum
of an absolutely continuous measure $\nu_c$
and a singular measure $\nu_s$,
both {\wrt} $\mu$:

\begin{theorem}\label{LEBESGUE_DECOMPOSITION_THEOREM}
Let $\mu$ be a positive $\sigma$-finite measure
and $\nu$ a complex measure on the measurable space $(X,\salgB)$.
Then there is a uniquely defined pair $(\nu_c,\nu_s)$ of complex measures
such that
\[
\nu = \nu_c + \nu_s, \quad
\nu_c \ll \mu, \quad
\nu_s \bot \mu.
\]
If the measure $\nu$ is positive,
then both measures $\nu_c$ and $\nu_s$ are positive measures, too.
\end{theorem}

%------------------------------------------------------------------------------%
%                                 BIBLIOGRAPHY                                 %
%------------------------------------------------------------------------------%
\begin{comment}
\vspace{8mm}
\begin{small}
\begin{thebibliography}{99}

\bibitem{El_2} J.\,Elstrodt:
\textit{Ma{\ss} -und Integrationstheorie}.
Springer Verlag, 2018.

\end{thebibliography}
\end{small}
\end{comment}

%\newpage
%\thispagestyle{empty}
\fancyhead[OC]{\small{\textsc{III. Measures on Locally Compact Spaces}}}
\fancyhead[EC]{\small{\textsc{III. Measures on Locally Compact Spaces}}}
\part{Measures on Locally Compact Spaces}
%-------------------------------------------------------------------------------
Abstract Measure theory also includes the study of measures
on topological spaces.
But these spaces must be supposed to be \textit{locally compact}
in order to obtain a satisfactory theory.

\subparagraph{}
Recall that the $\sigma$-algebra $[\cal{F}]$,
generated by a family of subsets $\cal{F} \subset 2^{X}$ of a base set $X$,
can be received as the intersection of all $\sigma$-algebras
containing this family $\cal{F}$.
If $(\topX,\topT)$ is a topological space,
the $\sigma$-algebra $[\topT]$ generated by the family $\topT$
of open sets is called the \textit{Borel $\sigma$-algebra} of $\topX$.
It is also generated by the family of closed sets of $\topX$.
Each set $B \xeof [\topT]$ is called \textit{Borel-measurable}
or a \textit{Borel set}.
Moreover,
if $\topX$ is Hausdorff and can be represented
as a union of mostly countable compact sets,
then $[\topT]$ is also generated by the family of compact subsets of $\topX$.

\subparagraph{}
The (topological) \textit{support} of a measure $\mu$ on a topological space
$\topX$ is a notion of where in the space $\topX$ the measure \textit{lives}.
It is defined to be the largest closed subset $\supp{\mu}$ of $\topX$,
for which every open neighbourhood of any point $x$ of $\supp{\mu}$
has positive measure.
A measure $\mu$ is called \textit{discrete},
if $\supp{\mu}$ consists of a mostly countable set of points,
i.e. if $\mu$ is concentrated to a sequence of single points
of the space $\topX$.

\subparagraph{}
Remember also
that bounded subsets of the Euclidean spaces $\Rd$ have a finite measure.
This property can be extended to topological spaces by the request,
that a measure should adopt finite values for compact sets.
Consequently we define a measure with this property as follows:
let $(\topX,\topT)$ be a topological space and $[\topT]$ its
Borel $\sigma$-algebra.
Then any measure $\mu$ is called a \textit{Borel measure} for $\topX$,
if $\mu(\compK) < \infty$ for every compact set $\compK \xeof [\topT]$.
For example,
the counting measure is a Borel measure on the set of integers
{\wrt} the discrete topology.
In addition,
a measure $\mu$ on the Borel sets of $\topX$
is called \textit{locally finite},
if for each point $x \xeof \topX$
there is an open neighbourhood $\openO_x$
such that $\mu(\openO_x) < \infty$.

\subparagraph{}
Each local finite measure is a Borel measure.
But if the topological space $\topX$ is a locally compact Hausdorff space,
then the inverse statement is also true,
that is,
each locally finite measure is a Borel measure.

\section{Regularity of Measures}
%------------------------------------------------------------------------------%
Let $\topX$ be a topological space and $\mu$ be a positive measure on its
Borel $\sigma$-algebra.
Then a Borel set $B \subseteq \topX$
is called \textit{inner regular} {\wrt} $\mu$, if
\[
\mu(B) \= \sup \,
\set{\mu(\compK)} {\compK \subset B, \, \compK \mbox{ compact}}.
\]
A measure $\mu$ is called \textit{inner regular},
if all Borel sets are inner regular.
Thus the inner regular measures on $\topX$ are already defined
by its values for compact subsets.
Each locally finite and inner regular measure on a topological space
$\topX$ is called a \textit{Radon measure}.

\subparagraph{}
On the other hand,
a measure $\mu$ on a topological space $\topX$
is called \textit{outer regular},
if for each Borel set $B \subseteq \topX$ the equality
\[
\mu(B) \= \inf \,
          \set{\mu(\openO)}{B \subset \openO, \, \openO \mbox{ open}}
\]
holds.
A measure is called \textit{regular},
if it is outer and inner regular.
For any finite and regular measure $\mu$,
for each Borel set $B \subset \topX$ and for each $\epsilon > 0$
there is a compact set $\compK_{\epsilon}$
and an open set $\openO_{\epsilon}$
such that
\[
\compK_{\epsilon} \subset B \subset \openO_{\epsilon}
\, \mbox{ and } \,
\mu(\openO_{\epsilon} \setminus \compK_{\epsilon}) < \epsilon.
\]
In other words,
for a Borel set $B$
the value $\mu(B)$ of such a measure $\mu$ can be approximated
by its values for compact subsets and open supersets of $B$.

\paragraph*{}
There are a couple of sufficient conditions on the topology in $\topX$,
which ensure
that Borel measures are regular.
For example,
this is true
if either $\topX$ is a compact metric space satisfying the
\textit{first countability axiom}
or $\topX$ is a \textit{polish space},
that is,  
$\topX$ is metrizable, separable and complete
(this is the assertion of \textsc{Ulam}\textit{'s theorem}).
If $\topX$ is a finite interval $[a,b]$,
this theorem implies
that every Borel measure on $[a,b]$ is Radon.

\subparagraph{}
The meaning of regular measures primarily relies on the fact
that so-called \textit{density theorems} can be quoted.
They ensure the possibility to approximate Lebesgue-integrable functions
by continuous ones.

\section{Haar Measures}
%------------------------------------------------------------------------------%
The existence of a translation invariant measure on a certain class of groups
is one of the highlights in abstract measure theory and
forms the foundation of a further mathematical area
called \textit{abstract harmonic analysis}.
 
\subparagraph{}
Remember
that a topological group is an algebraic one together with a topology $\topT$,
such that the group operations are continuous mappings.
A measure $\mu$ defined on the Borel sets of a topological group $\topG$
is called \textit{left invariant} respectively \textit{right invariant},
if for each element $g \xeof \topG$ and each Borel set
$B \xeof \cal{B}(\topG)$ we have
\[
\mu(B) \= \mu(g \circ B) \quad \mbox{rsp.} \quad
\mu(B) \= \mu(B \circ g).
\]
In addition,
a measure is called \textit{invariant},
if it is left \textit{and} right invariant.
Note
that if $\topG$ is an Abelian group,
the notions of left and right invariance of a measure fall together.

\subparagraph{}
Let $\topG$ be a locally compact Hausdorff space
and a topological group together.
Then there is a left invariant positive Radon measure
$\mu \: \cal{B}(\topG) \xto [0,\infty]$,
which is unique up to a constant $c > 0$.  
This measure is called a \textit{Haar measure} on $\topG$.

\paragraph{}
The following important examples of Haar measures are known:

\begin{itemize}[leftmargin=12pt]
\item[1.]
For the topological groups $(\R^d,+)$ the Haar measure is called the
\textit{Lebesgue-Borel measure} $\lambda^n$.

\item[2.]
For the complex unit circle $\Torus \subset \C$ the mapping
\[
\phi(f) \eqdef
\frac{1}{2\pi} \, \int_0^{2\pi} f(e^{it})\,dt, \quad f \in C(\C^{\times},\R)
\]
is a nonnegative invariant linear form,
and the associated invariant Radon measure is proportional
to the Lebesgue-Borel measure on $\Torus$.
\end{itemize}

\section{The Lebesgue-Borel Measure}
%------------------------------------------------------------------------------%
There are only few measurable spaces,
on which a measure can be defined in a very natural way.
An example of such measure spaces is any countable set,
whose $\sigma$-algebra is defined to be the power set of that space,
and a measure deserving the notion \textit{standard measure} for this space
is the counting measure $\mu_{\#}$,
that is,
the mapping assigning to each subset its cardinal number.
In particular,
for the set of integers
we get the measure space $(\Z,2^{\Z},\mu_{\#})$ this way.

\paragraph{}
Concerning the Euclidean spaces $\Rd$,
the situation is not quite so simple.
Here only those measures come into consideration,
which provide for such elementary sets like intervals,
rectangles or $d$-dimensional cubes
the associated lengths, areas or volumes.
The question arises
whether a positive measure exists on the Borel $\sigma$-algebra
$\Borel{{\R}^d}$ satisfying these properties.

\subparagraph{}
There is a fundamental theorem in measure theory,
known as the \textit{measure extension theorem},
which states
that in each space $\Rd$
there is exactly one positive measure $\LB^d$,
which adopts for each cube $W \subset \Rd$ with border lengths $l_1,\ldots,l_d$
the value
\[
\LB^d(W) \= l_1 \cdot l_2 \cdot \ldots \cdot l_{d-1} \cdot l_d,
\]
called the \textit{Lebesgue-Borel measure} on the spaces $\Rd$.

\subparagraph{}
It is not surprising 
that the measure $\LB^d$ is the Haar measure for the space $\Rd$.
Thus the Lebesgue-Borel measure have the property
to be invariant under Euclidean movements,
that is,
the equation $\LB^d(B) \eq \LB^d(\tau(B))$ holds
for each Borel set $B$ and
each orthogonal transformation $\tau \: \Rd \xto \Rd$,
in contrary to all other thinkable measures on space $\Rd$.

%------------------------------------------------------------------------------%
\section{Decomposition of Measures}                                            %
%------------------------------------------------------------------------------%
We turn back to Lebesgue's decomposition theorem, 
now considered for the measurable space $(\R,\Borel{\R})$
and the Lebesgue-Borel measure $\LB^1$.
In it general form it claims
that each positive measure $\nu$ can be represented in the form
\[
\nu = \nu_{ac} + \nu_d + \nu_{cs},
\]
$\nu_{ac} \ll \lambda$ being absolutely continuous,
$\nu_d$ being a measure 
which is concentrated on a possibly terminating sequence $(x_n)$
of points with $\nu_d(x_n) >0$,
and $\nu_{cs}$ being a \textit{continuous singular measure}.
A continuous singular measure is characterized by the fact
that it is concentrated on a set of measure $0$,
\textit{atomic} but not absolutely continuous w.r.t. $\lambda$.
An example of such a measure is the so-called \textit{Cantor distribution}.

\subparagraph{}
To each finite measure $\mu$ on $(\R,\Borel{\R})$ 
the monotonic increasing function
\[
F_{\mu}(a) \eqdef \mu(-\infty,a), \quad a \in \R
\]
is called the \textit{distribution function} of $\mu$.
Now, if $\mu \ll \lambda$ is absolutely continuous,
then the associated distribution function is called absolutely continuous, too.
In general,
each mapping $f \: [a,b] \xto \R$ is called \textit{absolutely continuous},
if there is a locally integrable function $g \: [a,b] \xto \R$,
such that
\[
f(x) \eqdef \int_a^x g(t) \, dt, \quad x \in [a,b].
\]
An absolutely continuous function has the property to be
$\lambda$-nearly everywhere differentiable,
and the derivative of $f$ is given by the function $g$.
One can characterize absolutely continuous functions as follows:
$f \: [a,b] \xto \R$ is \textit{absolutely continuous},
iff for each $\epsilon > 0$ there is a number $\delta > 0$
such that for each decomposition $\cal{Z}$ of $[a,b]$ with
grid size $F_{\cal{Z}} < \delta$
we have
$\sum_i \abs{ f(x_{i+1}) - f(x_i) } < \epsilon$.
The set of absolutely continuous functions on the interval $[a,b]$
forms a linear space denoted $\AC{[a,b]}$.

\section{Comparison of Lebesgue and Riemann Integration}
%------------------------------------------------------------------------------%
To the end,
let us have a look to the question
how the Lebesgue integral is comparable with the classical
Riemann integral for ordinary functions $f \: [a,b] \xto \R$.
For this purpose we will review in brief 
the setup and the basic
properties of the Riemannian integral.

\subparagraph{}
To define this concept for a function $f \: [a,b] \xto \R$,
two auxiliary concepts are required:

\begin{itemize}[leftmargin=12pt]
\item[1.]
A \textit{decomposition} $P$ of $[a,b]$ is a sequence of points
\[
\{x_0 := a,x_1,\ldots,x_{n-1},x_n := b \},
\quad x_{i-1} < x_i, \quad i=1,\ldots,n,
\]
and the length $\snorm{P} $ of the largest subinterval
is called the \textit{grid} of $P$,
that is
\[
\snorm{P} \eqdef \max_{1 \leq i \leq n} \, (x_i-x_{i-1}).
\]
\item[2.]
A \textit{tagged} decomposition $P(t_0,\ldots,t_n)$
of the interval $[a,b]$
is a decomposition $P$ (in the ordinary sense) together with a 
sequence of points $t_0,\ldots,t_{n-1}$,
where each point $t_i$ is contained in the subinterval $[x_i,x_{i+1}]$.
\end{itemize}

\subparagraph{}
Now, 
for a tagged decomposition $P(t_0,\ldots,t_n)$ 
the \textit{Riemannian sum} $R(f,P)$
for the function $f\:[a,b] \xto \R$ is defined by
\[
R(f,P) \eqdef \sum_{i=1}^n f(t_i) \cdot (x_i - x_{i-1}),
\]
and the \textit{Riemann integral} of $f$ is received as
possible limit value of Riemannian sums $R(f,P)$ by decreasing the
grid to zero.
The function $f$ is called \textit{Riemann integrable} on $[a,b]$
(briefly \textit{R-integrable}) to the value $R_f$,
if for every $\epsilon > 0$ there is a number $\delta > 0$,
such that
$\abs{ R_f - R(f,P)} < \epsilon$
for every tagged decomposition $P(t_0,\ldots,t_n)$
with grid $\abs{P} < \delta$,
and the magnitude $R_f$ is said to be the \textit{Riemann integral} of $f$
in this case,
usually denoted
\[
\int_a^b f \, dx \eqdef R_f.
\]
It is easy to verify
that the set $R(a,b)$ of Riemann integrable functions
on the interval $[a,b]$ forms a linear space,
so Riemann integration appears as a linear functional $R \: \R(a,b) \xto \R$.

\subparagraph{}
Notice also that
each continuous function,
each monotonic function and
each function of bounded variation 
(being the difference of two monotonic increasing functions)
is a member of the space $R(a,b)$.

\paragraph{}
The following theorem goes back to \textsc{Lebesgue}
and provides a connection between both integration methods
for functions defined on a compact interval $[a,b] \subset \R$:

\begin{theorem}
Each R-integrable function $f\:[a,b] \xto \R$ is
Lebesgue integrable {\wrt} the Lebesgue-Borel measure $\LB := \LB^1$).
In addition,
for these functions both kinds of integrals coincide, {\ie}
\[
\int_a^b f(x) \, dx \= \int_{[a,b]} f \, d\LB.
\]
\end{theorem}

\subparagraph{}
Due to this theorem it is possible to compute the Lebesgue integral
by applying the known integration rules for
Riemann integration like
the \textit{product rule},
the \textit{substitution rule} and
the \textit{formula of partial integration}.

%------------------------------------------------------------------------------%
%                                 BIBLIOGRAPHY                                 %
%------------------------------------------------------------------------------%
\begin{comment}
\vspace{8mm}
\begin{small}
\begin{thebibliography}{99}

\bibitem{El_3} J.\,Elstrodt:
\textit{Ma{\ss} -und Integrationstheorie}.
Springer Verlag, 2018.

\end{thebibliography}
\end{small}
\end{comment}


%------------------------------------------------------------------------------%
%                                 BIBLIOGRAPHY                                 %
%------------------------------------------------------------------------------%
\vspace{8mm}
\begin{thebibliography}{99}

\bibitem{El_1} J.\,Elstrodt:
\textit{Ma{\ss} -und Integrationstheorie}.
Springer Verlag, 2018.

\end{thebibliography}

%------------------------------------------------------------------------------%
%                               END OF DOCUMENT                                %
%------------------------------------------------------------------------------%
\end{document}
